\documentclass[12pt]{article}

\usepackage{sbc-template}
\usepackage[utf8]{inputenc}
\usepackage[T1]{fontenc}
\usepackage[brazil]{babel}
\usepackage{graphicx,url}
\usepackage{booktabs}
\usepackage{multirow}
\usepackage{xcolor}
\usepackage{tikz}
\usetikzlibrary{positioning,fit,arrows.meta,backgrounds}

% Marca visível de pendência; trocar a definição por \newcommand{\pendente}[1]{} antes de submeter.
\newcommand{\pendente}[1]{{\color{red}\textbf{[PENDENTE: #1]}}}
\newcommand{\kmp}{\textsc{KMP}}

\sloppy

\title{KMP: Distribuição Cooperativa e Autenticada de Pacotes\\
       para Consoles Portáteis Linux de Baixo Custo}

\author{Weslei Ferreira Santos\inst{1}}

\address{\pendente{instituição, departamento, cidade -- UF}
  \email{\pendente{e-mail institucional}}
}

\begin{document}

\maketitle

\begin{abstract}
  Linux handheld consoles with 1~GB of RAM receive community game ports
  through centralized HTTP downloads, which concentrate cost on a few hosts and
  ignore the idle capacity of the devices themselves. This paper presents
  KMP, a tool that distributes packages from an Ed25519-signed catalog over
  BitTorrent, with a standalone static daemon for aarch64 consoles, local
  discovery, IPv6 connectivity between home networks backed by a public seed
  node, a desktop publisher and battery-aware community seeding. In an A/B test on a
  1~GB console, the current engine kept transfer speed while seeding with
  27~MB of resident memory, against 271~MB in the previous version.
\end{abstract}

\begin{resumo}
  Consoles portáteis Linux com 1~GB de RAM recebem ports de jogos da
  comunidade por downloads HTTP centralizados, que concentram o custo em poucos
  servidores e ignoram a capacidade ociosa dos próprios aparelhos. Este artigo
  apresenta o KMP, uma ferramenta que distribui pacotes de um catálogo assinado
  com Ed25519 por BitTorrent, com um daemon estático e independente para
  consoles aarch64, descoberta local, conectividade IPv6 entre redes
  domésticas apoiada por um nó público, um publicador para desktop e semeação
  comunitária ciente da bateria. Num teste A/B em um console de 1~GB, o motor atual manteve a
  velocidade de transferência semeando com 27~MB de memória residente, contra
  271~MB na versão anterior.
\end{resumo}

% ---------------------------------------------------------------------------
\section{Introdução}
\label{sec:intro}

Consoles portáteis baseados em Linux, como os modelos R36H e R46H, usam
processadores ARM de 64 bits, 1~GB de RAM e cartões SD, e rodam firmwares
comunitários como ArkOS, dArkOS e ROCKNIX. O ecossistema
PortMaster~\cite{portmaster} disponibiliza para esses aparelhos centenas de
ports, pacotes que vão de poucos megabytes a vários gigabytes. A distribuição é
feita por downloads HTTP de servidores e repositórios centrais: cada
instalação consome banda de quem hospeda o arquivo, e os consoles que já
possuem o pacote não ajudam os demais.

Distribuição cooperativa por BitTorrent~\cite{cohen2003incentives,bep3} é uma
resposta conhecida para esse custo, e sistemas como
Dragonfly~\cite{dragonfly}, Kraken~\cite{kraken} e
PeerSync~\cite{deng2026peersync} a aplicam a imagens de contêineres em
datacenters e na borda. O cenário dos consoles, porém, combina restrições que
esses sistemas não têm como premissa:
(i)~os participantes são aparelhos de 1~GB de RAM, alimentados por bateria,
com firmwares imutáveis em que não é possível instalar dependências no
sistema;
(ii)~estão espalhados por redes domésticas heterogêneas, frequentemente atrás
de NAT e sem portas de entrada; e
(iii)~não há operador central: os pacotes vêm de autores da comunidade e
precisam ser autenticados antes de instalados.

Este artigo investiga a pergunta: \emph{quais limitações operacionais
dificultam a distribuição cooperativa de arquivos em dispositivos embarcados
de baixo custo espalhados por redes domésticas, e em que medida uma
ferramenta integrada consegue superá-las?} Como resposta, apresentamos o
\kmp{}, que reúne em uma única ferramenta a publicação, a autenticação, a
distribuição, a instalação e a semeação comunitária de pacotes. As
contribuições são:

\begin{itemize}
  \item um daemon BitTorrent estático para aarch64, independente da interface
        e controlado por JSON-RPC, com descoberta local, IPv6 e hole punching,
        cujo consumo de recursos foi medido em um console de 1~GB
        (Seção~\ref{sec:avaliacao});
  \item um catálogo assinado com Ed25519, distribuído pelo próprio swarm, com
        proteção contra rollback e múltiplos espelhos sem chave privada
        (Seção~\ref{sec:arquitetura});
  \item a medição, em redes domésticas reais, de por onde chega cada byte: na
        mesma rede, os próprios consoles entregam a maior parte dos pacotes; em
        redes diferentes, o nó público garante a entrega imediata e os consoles
        assumem por IPv6 e uTP quando a descoberta se propaga
        (Seção~\ref{sec:avaliacao}).
\end{itemize}

Nenhum desses mecanismos é inédito isoladamente. A contribuição está em
integrá-los e em operá-los sobre hardware restrito, com avaliação
experimental. A Seção~\ref{sec:relacionados} discute os trabalhos
relacionados; a Seção~\ref{sec:arquitetura} descreve a arquitetura; a
Seção~\ref{sec:avaliacao} apresenta a avaliação; a
Seção~\ref{sec:artefatos} lista os artefatos e a demonstração planejada; e a
Seção~\ref{sec:conclusao} conclui.

% ---------------------------------------------------------------------------
\section{Trabalhos relacionados}
\label{sec:relacionados}

\textbf{Clientes BitTorrent de uso geral.} Transmission~\cite{transmission} e
aria2~\cite{aria2} implementam o protocolo com DHT~\cite{bep5},
PEX~\cite{bep11} e descoberta local~\cite{bep14}, e oferecem interfaces RPC.
São ferramentas de transferência: não decidem quais arquivos são legítimos,
não instalam pacotes e não foram desenhadas para conviver com uma interface
gráfica na mesma memória de um console. O \kmp{} usou o aria2 como motor nas
primeiras versões e o mantém como alternativa de reserva.

\textbf{Distribuição P2P em datacenters e na borda.} O
Dragonfly~\cite{dragonfly}, projeto graduado da CNCF, e o
Kraken~\cite{kraken} aceleram a distribuição de imagens de contêineres entre
nós de um cluster, com componentes de agendamento e rastreamento. O
PeerSync~\cite{deng2026peersync} leva a ideia para a borda, com eleição
automática de tracker, recuperação de falhas e liberação cooperativa de
cache, e relata ganhos de 1,79$\times$ sobre o Dragonfly e 1,28$\times$ sobre
o Kraken. Esses sistemas pressupõem nós sob um mesmo domínio administrativo e
integração com runtimes de contêineres; o \kmp{} atende aparelhos de usuários
finais, em redes que ninguém administra.

\textbf{Atualização segura e travessia de NAT.} O TUF~\cite{samuel2010tuf}
sistematiza ameaças a sistemas de atualização, como rollback e chaves
comprometidas; o catálogo do \kmp{} adota parte dessas defesas, com sequência
monotônica e expiração. A travessia de NAT por hole
punching~\cite{ford2005nat,rfc5128} é padronizada no BitTorrent pela
BEP~55~\cite{bep55}.

A Tabela~\ref{tab:comparacao} resume as diferenças segundo a documentação
consultada. A ausência de marca não significa ausência da funcionalidade, e
sim que ela não foi confirmada nas fontes.

\begin{table}[ht]
\centering
\caption{Características segundo a documentação de cada ferramenta
($\bullet$ documentado; ? não confirmado; -- fora do escopo).}
\label{tab:comparacao}
\footnotesize
\setlength{\tabcolsep}{4pt}
\begin{tabular}{@{}lccccc@{}}
\toprule
Característica & \kmp{} & Transmission & aria2 & Dragonfly & PeerSync \\
\midrule
Alvo principal & consoles & desktop & desktop & cluster & borda \\
DHT, PEX e descoberta local & $\bullet$ & $\bullet$ & $\bullet$ & ? & ? \\
Hole punching (BEP 55) & $\bullet$ & ? & ? & ? & ? \\
Catálogo assinado com anti-rollback & $\bullet$ & -- & -- & ? & ? \\
Instalação dos pacotes & $\bullet$ & -- & -- & $\bullet$ & $\bullet$ \\
Executável estático, sem dependências & $\bullet$ & ? & ? & ? & ? \\
Avaliado em aparelho de 1~GB & $\bullet$ & ? & ? & ? & ? \\
\bottomrule
\end{tabular}
\end{table}

\pendente{conferir cada ``?'' na documentação antes da submissão.}

% ---------------------------------------------------------------------------
\section{Arquitetura e funcionalidades}
\label{sec:arquitetura}

O \kmp{} tem cinco componentes (Figura~\ref{fig:arquitetura}): o
\emph{Publisher}, o \emph{catálogo assinado}, o daemon
\emph{kmp-torrentd}, o \emph{Client} do console e um \emph{nó público} de
semeação. Somados, o daemon, o Client e os scripts do console têm cerca de
7.800 linhas de Go, Lua, Python e shell.

\begin{figure}[ht]
\centering
\begin{tikzpicture}[
    font=\footnotesize,
    box/.style={draw, rounded corners=2pt, align=center, minimum height=8mm, inner sep=3pt},
    group/.style={draw, dashed, rounded corners=3pt, inner sep=5pt},
    arr/.style={-{Stealth[length=2mm]}, thick}]
  \node[box] (pub) {Publisher\\(Electron)};
  \node[box, right=16mm of pub] (ci) {Repositório do catálogo\\validação no CI};
  \node[box, right=12mm of ci] (head) {\texttt{head.json} + assinatura\\espelhos HTTPS};
  \node[box, below=14mm of ci] (swarm) {Swarm BitTorrent\\DHT $\cdot$ PEX $\cdot$ LPD $\cdot$ BEP 55};
  \node[box, left=10mm of swarm] (node) {Nó público\\(semeador)};
  \node[box, below=12mm of swarm, xshift=-22mm] (ui) {Interface\\(LÖVE)};
  \node[box, right=8mm of ui] (daemon) {kmp-torrentd\\(Go, estático)};
  \node[box, right=8mm of daemon] (comm) {Nó comunitário\\e bateria};
  \node[group, fit=(ui)(daemon)(comm), label={[font=\footnotesize]below:Client no console aarch64}] (client) {};
  \draw[arr] (pub) -- node[above]{proposta} (ci);
  \draw[arr] (ci) -- node[above]{assina} (head);
  \draw[arr] (ci) -- node[right]{torrents} (swarm);
  \draw[arr, <->] (node) -- (swarm);
  \draw[arr, <->] (daemon) -- (swarm);
  \draw[arr] (head) |- node[pos=.25, right]{verifica} (client.east);
  \draw[arr, <->] (ui) -- node[above]{RPC} (daemon);
  \draw[arr, <->] (comm) -- (daemon);
\end{tikzpicture}
\caption{Arquitetura do KMP: a publicação passa pelo repositório do catálogo;
os pacotes circulam pelo swarm, do qual os consoles participam.}
\label{fig:arquitetura}
\end{figure}

\textbf{Publisher.} Aplicação Electron para Linux, Windows e macOS. O autor
seleciona o port; o Publisher monta o pacote ZIP, calcula o SHA-256, cria o
torrent, abre a proposta como pull request no repositório do catálogo e
semeia o pacote a partir do computador do autor. O Publisher também tem
atualização automática verificada por SHA-256.

\textbf{Catálogo assinado.} Cada publicação gera um snapshot JSON imutável
com número de sequência crescente, assinado com Ed25519~\cite{rfc8032} sobre
os bytes exatos do arquivo. Um documento pequeno, o \texttt{head.json},
também assinado, aponta para o snapshot pelo magnet, pelo tamanho e pelo
SHA-256, e tem data de expiração. O mesmo par é copiado para vários
espelhos, que não possuem a chave privada. O console busca o head em todos
os espelhos, rejeita assinaturas inválidas, heads expirados e sequências
menores que a maior já aceita, e baixa o snapshot pelo próprio swarm. Se
todas as fontes falharem, continua com o último snapshot válido. Antes de
assinar, o CI baixa cada pacote proposto pelo swarm e confere tamanho e
hash.

\textbf{kmp-torrentd.} Daemon em Go sobre a biblioteca
anacrolix/torrent~\cite{anacrolix}, compilado de forma estática para
aarch64. Por não depender de bibliotecas do sistema, roda em firmwares
imutáveis sem modificar a partição do sistema. Implementa TCP,
uTP~\cite{bep29}, DHT IPv4 e IPv6, PEX, descoberta local por multicast
(BEP~14) e hole punching (BEP~55), e expõe uma API JSON-RPC compatível com o
subconjunto do aria2 usado pelo Client, o que permitiu trocar de motor sem
reescrever a interface. O daemon é independente da interface: sobrevive ao
fechamento do Client, restaura a sessão e continua semeando. Três decisões
vêm da operação em 1~GB de RAM:
(i)~os arquivos são lidos e escritos sem mapeamento em memória
(Seção~\ref{sec:recursos});
(ii)~em aparelhos de até 1,5~GB, recebe um teto suave de 192~MiB para o heap
do Go; e
(iii)~o download não tem limitador de taxa por padrão, porque, na biblioteca,
um limitador esgotado impede conexões novas, e o console deixava de alcançar
o nó público.
Cada console sorteia uma porta de escuta, para que vários aparelhos na mesma
rede possam ser mapeados por UPnP/NAT-PMP.

\textbf{Client.} Interface em LÖVE~\cite{love2d} no formato de loja, com
abas de início, instalados e atualizações, controlada pelos botões do
console. A interface limita a taxa de quadros a 30 por segundo em uso e a 8
quando parada. O Client baixa e confere cada pacote antes de instalá-lo e
pode se atualizar pelo próprio catálogo, numa transação com journal que
sobrevive a quedas de energia. Um \emph{nó comunitário} opcional baixa e
semeia pacotes do catálogo dentro de uma cota de armazenamento, com no máximo
dois downloads simultâneos, e pausa a semeação com a bateria abaixo de 30\%.

\textbf{Nó público.} Uma instância gratuita de nuvem (1~vCPU, 1~GB) roda o
mesmo daemon e semeia todo o catálogo. É anunciado no catálogo e adicionado
explicitamente a cada download, garantindo uma fonte alcançável quando
nenhum console está ligado, e coordena o hole punching entre consoles atrás
de NAT.

% ---------------------------------------------------------------------------
\section{Avaliação}
\label{sec:avaliacao}

A avaliação responde a três questões:
\textbf{Q1:} quanto o \kmp{} consome em um console de 1~GB, e as otimizações
de memória custaram velocidade de transferência?
\textbf{Q2:} consoles na mesma rede e em redes diferentes conseguem semear
entre si sem configuração manual?
\textbf{Q3:} como o \kmp{} se compara a um cliente BitTorrent de uso geral
nas mesmas condições?

\subsection{Consumo de recursos (Q1)}
\label{sec:recursos}

\textbf{Metodologia.} Comparamos os motores das versões 0.2.56 (com
mapeamento de arquivos em memória) e 0.2.60 (leitura comum) no mesmo
console, com processador do R36H/R46H e 895~MB de RAM utilizáveis, com as
mesmas opções (sem limite de banda e sem descoberta local) e três rodadas por
versão em ordem alternada. Em cada rodada, o console baixa 256~MiB de dados
aleatórios só de um PC pela rede Wi-Fi doméstica, fica 45~s parado e envia o
pacote inteiro de volta. CPU e memória são lidas de \texttt{/proc} a cada
segundo, no próprio console. A interface foi medida em três janelas de 60~s
com a loja aberta e parada, trocando apenas o código da interface.

\begin{table}[ht]
\centering
\caption{Teste A/B no console de 1~GB: mediana e, entre parênteses, a menor e
a maior das três rodadas.}
\label{tab:ab}
\small
\begin{tabular}{@{}lll@{}}
\toprule
Métrica & 0.2.56 & 0.2.60 \\
\midrule
Download (MiB/s) & 2,19 (2,00--2,96) & 2,89 (2,14--3,06) \\
Envio (MiB/s) & 4,20 (3,96--4,25) & 4,25 (4,12--4,25) \\
CPU do motor baixando (\% de um núcleo) & 55 (55--59) & 55 (54--60) \\
CPU do motor enviando (\%) & 19 (16--37) & 23 (22--45) \\
CPU do motor parado (\%) & 2,0 (0,8--2,9) & 2,6 (1,6--2,9) \\
Memória do motor, pico (RSS, MB) & 271 (263--281) & 27 (27--30) \\
\quad memória anônima (MB) & 16 (15--17) & 15 (15--16) \\
\quad pacote mapeado (MB) & 259 (247--265) & -- \\
RAM disponível no pior momento (MB) & 322 (322--386) & 328 (316--350) \\
CPU da interface parada (\%) & 34,9 (34,8--34,9) & 5,1 (5,05--5,15) \\
\bottomrule
\end{tabular}
\end{table}

\textbf{Resultados.} A Tabela~\ref{tab:ab} mostra que a velocidade de envio
foi a mesma e que as faixas de download se sobrepõem: a diferença de
mediana vem da variação do Wi-Fi, não de ganho. A CPU do motor também não
mudou. A memória residente do motor caiu 90\% durante a semeação, mas a
memória anônima, que o sistema não pode descartar, é a mesma nas duas
versões: a diferença inteira eram páginas do pacote mapeado, que o kernel
consegue liberar sob pressão, e por isso a RAM disponível ficou
praticamente igual. O ganho prático da leitura comum é o motor não reter os
pacotes que semeia e não ser encerrado por \texttt{SIGBUS} quando um arquivo
semeado é trocado ou o cartão é removido. O maior ganho visível é o da
interface, cuja CPU parada caiu 85\% com o limite de quadros.

\subsection{Conectividade entre consoles (Q2)}

\textbf{Metodologia.} Para saber por onde chegou cada byte, o daemon passou a
acumular, por torrent, os bytes recebidos e enviados por caminho, isto é, por
endereço e porta remotos, origem do peer (tracker, DHT, PEX, descoberta local,
peer passado pelo Client ou conexão de entrada) e transporte (TCP ou uTP sobre
IPv4 ou IPv6), incluindo as conexões já encerradas. Um console ficou na rede
doméstica, semeando o catálogo pelo nó comunitário, junto com o Publisher num
PC da mesma rede e o nó público. O outro console baixou pelo próprio Client,
com a opção ``Só P2P'' (sem espelhos HTTP) e o nó comunitário desligado, três
pacotes do catálogo: primeiro na rede doméstica e depois conectado à rede
móvel de um celular, com IPv4 atrás de CGNAT e IPv6 público. A rede doméstica
não aceita conexões IPv4 de entrada, mas aceita IPv6. Um registrador no
próprio console gravou os totais a cada 2~s.

\begin{table}[ht]
\centering
\caption{Origem dos bytes recebidos pelo console em cada download: fração
entregue pelo nó público e pelos demais aparelhos (o outro console e o PC).
Pacotes: Asphalt~5 (108~MiB), Riptide GP (44~MiB) e Celeste~64 (31~MiB).}
\label{tab:conectividade}
\footnotesize
\setlength{\tabcolsep}{4pt}
\begin{tabular}{@{}llrrrrl@{}}
\toprule
Rede & Pacote & Min.\ após a troca & Duração (s) & Nó & Outros & Caminho dos outros \\
\midrule
\multirow{3}{*}{doméstica} & Asphalt 5 & -- & 48 & 29\% & 71\% & IPv4, TCP, LPD \\
 & Riptide GP & -- & 26 & 9\% & 91\% & IPv4, TCP, PEX e DHT \\
 & Celeste 64 & -- & 18 & 25\% & 75\% & IPv4, TCP, LPD \\
\midrule
\multirow{8}{*}{móvel} & Asphalt 5 & 0,8 & 74 & 100\% & 0\% & -- \\
 & Celeste 64 & 2,3 & 19 & 100\% & 0\% & -- \\
 & Riptide GP & 3,1 & 25 & 100\% & 0\% & -- \\
 & Asphalt 5 & 7,3 & 73 & 100\% & 0\% & -- \\
 & Celeste 64 & 9,2 & 22 & 100\% & 0\% & -- \\
 & Riptide GP & 9,9 & 26 & 67\% & 31\% & IPv6 e IPv4, uTP, DHT \\
 & Asphalt 5 & 11,3 & 49 & 27\% & 73\% & IPv6 e IPv4, uTP, DHT \\
 & Celeste 64 & 12,5 & 16 & 35\% & 64\% & IPv6, uTP, DHT \\
\bottomrule
\end{tabular}
\end{table}

\textbf{Resultados.} Na rede doméstica, o outro console e o PC entregaram de
71\% a 91\% de cada pacote, encontrados pela descoberta local, por PEX e pela
DHT, em IPv4 e IPv6 (Tabela~\ref{tab:conectividade}). Na rede móvel, o padrão
depende do tempo desde a troca de rede. Nos downloads iniciados até cerca de
9~minutos depois, o nó público entregou praticamente 100\%: o console conhecia
os outros aparelhos pelos endereços IPv4 da rede doméstica, mas não conseguia
alcançá-los, e o daemon registrou 19 peers indiscáveis sem hole punching e uma
única tentativa de hole punching, sem sucesso. A partir de 10~minutos, os
aparelhos da rede doméstica passaram a entregar de 31\% a 73\% de cada pacote,
quase tudo por uTP e descobertos pela DHT, majoritariamente por IPv6, e o
Asphalt~5 caiu de 73~s para 49~s. Esse intervalo é coerente com o tempo de
propagação do novo endereço do console pela DHT, já que o nó público, só com
IPv4, não repassa endereços IPv6 por PEX. Numa rodada adicional com o nó
comunitário do console ligado, os outros aparelhos também só entraram no meio
dos downloads, o que indica que manter o console nos swarms não elimina esse
intervalo.

Os resultados mostram as duas funções da infraestrutura: o nó público garante a
entrega logo após uma mudança de rede, e a semeação entre consoles assume a
maior parte do tráfego quando a descoberta se propaga. Além disso, um pacote de
2,96~GiB foi validado pelo CI baixando exclusivamente do swarm em 21,5~min. O
experimento também revelou dois defeitos, a corrigir na próxima versão: a
descoberta local não se recuperava após a troca de rede, e o gerenciador de
ports dos firmwares encerra todos os processos Python ao sair de qualquer jogo,
inclusive os do nó comunitário.

\subsection{Comparação com um cliente de uso geral (Q3)}

\textbf{Metodologia.} Repetimos o roteiro da Seção~\ref{sec:recursos} com o
\texttt{transmission-daemon} 4.0.6~\cite{transmission} no lugar do motor
antigo, em outro console do mesmo modelo. Como o Transmission não aceita
peers pela API, os dois motores, o semeador e o receptor no PC se encontram
pelo mesmo tracker HTTP na rede local. O Transmission foi obtido dos pacotes
do Alpine Linux e roda com as próprias bibliotecas musl, sem alterar o
firmware, e usa a configuração padrão, exceto pelo que precisa ser igual ao
\kmp{}: sem limites de banda e de proporção, sem descoberta local e com DHT,
PEX, uTP e mapeamento de porta ligados. Separamos o tempo até o primeiro byte
da velocidade medida a partir dele.

\begin{table}[ht]
\centering
\caption{kmp-torrentd 0.2.8 contra Transmission 4.0.6 no console de 1~GB:
mediana e, entre parênteses, a menor e a maior das três rodadas.}
\label{tab:transmission}
\small
\begin{tabular}{@{}lll@{}}
\toprule
Métrica & kmp-torrentd & Transmission \\
\midrule
Tempo até o primeiro byte (s) & 0,4 (0,4--0,5) & 10,6 (10,4--10,9) \\
Download após o primeiro byte (MiB/s) & 3,66 (2,81--4,64) & 2,80 (2,69--2,98) \\
Envio após o primeiro byte (MiB/s) & 4,63 (4,24--4,80) & 4,35 (4,22--4,57) \\
CPU por GiB baixado (s) & 146 (137--148) & 200 (186--207) \\
CPU por GiB enviado (s) & 50 (47--86) & 63 (61--65) \\
CPU parado (\% de um núcleo) & 3,1 (2,6--3,3) & 0,7 (0,7--0,8) \\
Memória, pico (RSS, MB) & 35 (35--36) & 29 (29--32) \\
\quad memória anônima (MB) & 19 (19--20) & 19 (19--23) \\
Memória parado, ao fim (RSS, MB) & 27 (26--30) & 11 (11--11) \\
Consultas à API sem resposta em 3~s & 0 & 1--3 por rodada \\
\bottomrule
\end{tabular}
\end{table}

\textbf{Resultados.} A Tabela~\ref{tab:transmission} mostra três vantagens do
\kmp{} com faixas que não se sobrepõem: começa a transferir em menos de meio
segundo, contra cerca de 10~s do Transmission; gasta 27\% menos CPU por GiB
baixado; e sua API respondeu a todas as consultas, enquanto a do Transmission
deixou de responder em até 3~s de uma a três vezes por rodada, sempre no início
da transferência. A velocidade do \kmp{} foi maior na mediana, mas as faixas
se sobrepõem e três rodadas não bastam para afirmar vantagem. O Transmission é
mais econômico parado: usa 0,7\% de CPU contra 3,1\% e termina com 11~MB de
memória residente contra 27~MB; a memória anônima dos dois é a mesma. Não
investigamos a causa do atraso inicial do Transmission, que pode depender de
ajustes fora da configuração padrão.

Em suma, nenhum dos dois domina em todas as métricas. O kmp-torrentd é mais
eficiente e responsivo durante a transferência, que é o momento crítico num
console que divide a CPU com a interface, e o Transmission gasta menos com o
aparelho parado. Esse custo em repouso é um ponto a melhorar no \kmp{}.

% ---------------------------------------------------------------------------
\section{Artefatos, documentação e demonstração}
\label{sec:artefatos}

\textbf{Código-fonte e licença.} O Client e o daemon estão em
\url{https://github.com/FerreiraTechLab/KMP} e o Publisher em
\url{https://github.com/FerreiraTechLab/KMP-Publisher}, ambos sob a
GPL-3.0; a cópia modificada da anacrolix/torrent mantém a MPL-2.0. O
catálogo (MIT) está em \url{https://github.com/FerreiraTechLab/KMP-catalog}
e os binários (Apache-2.0) em
\url{https://github.com/FerreiraTechLab/KMP-releases}.

\textbf{Documentação.} O manual de instalação, uso e arquitetura está em
\url{https://kpm-ferreiratechlab.web.app}. As medições da
Seção~\ref{sec:avaliacao}, com dados brutos, scripts e instruções de
reprodução, estão em \url{https://github.com/FerreiraTechLab/KMP-Artigo}.

\textbf{Demonstração planejada.} Dois consoles portáteis, um notebook e um
celular como roteador móvel. Roteiro: (1)~publicar um pacote pelo Publisher;
(2)~mostrar o catálogo assinado chegando aos consoles; (3)~baixar o pacote em
um console a partir do outro, na mesma rede; (4)~mudar um console para a
rede móvel e mostrar a semeação entre redes pelo nó público. Requisitos:
uma tomada, rede Wi-Fi com acesso à Internet e uma mesa.

% ---------------------------------------------------------------------------
\section{Conclusão}
\label{sec:conclusao}

O \kmp{} mostra que é possível operar distribuição cooperativa e autenticada
de pacotes em consoles de 1~GB espalhados por redes domésticas, com um
daemon independente, consumo de recursos medido e semeação entre redes
diferentes. Comparado ao Transmission no mesmo console, o daemon do \kmp{}
começou a transferir 10~s antes e gastou 27\% menos CPU por GiB baixado, mas
consome mais em repouso. Os trabalhos futuros incluem reduzir esse custo em
repouso, um nó público com IPv6, que poderia repassar endereços IPv6 e encurtar
o intervalo até os consoles se encontrarem em redes diferentes, nós públicos
redundantes operados pela comunidade e a integração com o gerenciamento de um
segundo cartão SD.

\bibliographystyle{sbc}
\bibliography{referencias}

\end{document}
